\section{Conclusion}
\label{sec:conclusion}

We studied teacher-guided knowledge distillation for the Regression Tsetlin Machine by
letting a frozen teacher steer the student's clause-level Type~I/II feedback, and asked
not only whether it helps but when, where, and why. Across a family of mechanisms and a
controlled masked-distillation benchmark, we found that every mechanism is a
\emph{sharpener} that scales the teacher's signal without rescuing a bad teacher; that
teacher accuracy does not predict benefit; that a 70\%-data guided student can reach and
occasionally exceed a full-data model; and that data spacing/density---not teacher
accuracy---governs distillation benefit. These consolidate into a single account:
distillation is a predominantly global reshaping of the learned function that helps when
a dataset's error is not spacing-governed and fails when it is, with a good teacher
additionally repairing the specific regions whose training support was removed. The
transparency of the TM's clause feedback was essential to reaching, and to explaining,
each of these conclusions.

\paragraph{Open threads.}
\begin{itemize}
  \item Distinguishing input-space spacing from joint $(x,y)$ spacing as the operative
    variable.
  \item Firming the per-dataset density$\leftrightarrow$benefit correlation for modes
    other than \texttt{force\_ii\_tmag}, where it is currently weaker.
  \item Testing the account under non-structured (e.g.\ random or adversarial) scarcity
    patterns.
\end{itemize}
